% Theory of the measurements. Generic: uses only macros from results.tex.

\subsection{Ideal quantizer}
An $N$-bit ADC maps its input onto $2^N$ codes. Code $k$ is output for inputs between the
transition levels $T_k$ and $T_{k+1}$; the ideal levels are equally spaced by one LSB,
$Q = (T_{2^N-1} - T_1)/(2^N - 2)$. For an input that is busy enough, the quantization error is
uniform in $[-Q/2, Q/2]$, with power $Q^2/12$. A full-scale sine has power
$(2^N Q)^2/8$, so an ideal ADC has
\begin{equation}
\text{SNR}_{ideal} = 10\log_{10}\frac{(2^N Q)^2/8}{Q^2/12} = 6.02\,N + 1.76\ \text{dB},
\end{equation}
\SI{31.9}{dB} for $N = \AdcBits$. A real ADC adds nonlinearity, noise and dynamic errors to the
quantization error; the effective number of bits is the resolution of an ideal ADC with the
same signal to noise and distortion ratio:
\begin{equation}
\text{ENOB} = \frac{\text{SINAD} - 1.76}{6.02}.
\end{equation}
The ENOB in this report is referred to the amplitude of the applied sine, which is smaller than
full scale (to avoid clipping). An ideal \AdcBits-bit quantizer driven by the same amplitude
reaches a little less than \AdcBits{} bits, about 4.7 bits for a sine of 27 LSB peak-to-peak.

\subsection{Static linearity: the histogram test}
The transition levels are measured with the code density (histogram) test of IEEE~1241. The input
is a triangle that overdrives both ends of the range, so the end codes $0$ and $2^N-1$ appear and
every inner code is crossed at a constant slope. With a linear ramp from $V_{min}$ to $V_{max}$ the
voltage is uniformly distributed in time, and the number of samples $H_k$ that land in code $k$ is
proportional to its width. With $M$ samples in total and the cumulative histogram
$C_k = \sum_{j \le k} H_j$, the transition levels are
\begin{equation}
T_k = V_{min} + (V_{max} - V_{min})\,\frac{C_{k-1}}{M}, \qquad k = 1 \ldots 2^N - 1.
\end{equation}
From them, with the end-point LSB $Q$ defined above,
\begin{equation}
\text{DNL}_k = \frac{T_{k+1} - T_k}{Q} - 1, \qquad
\text{INL}_k = \frac{T_k - \left(T_1 + (k-1)\,Q\right)}{Q}.
\end{equation}
DNL is the error of the width of each code (a code with $\text{DNL} = -1$ is missing); INL is the
deviation of each transition from the straight line through the first and last transitions, so
offset and gain errors are removed.

Three details make the test accurate:
\begin{itemize}
  \item \textbf{Whole periods.} The capture of $M$ samples must hold an integer number of ramp
        periods; otherwise the leftover fraction of a period lands on different codes from run to
        run and biases the histogram. The ramp frequency is $f = 10\,f_s/M$ (10 periods).
  \item \textbf{Ramp ends.} $V_{min}$ and $V_{max}$ come from the scope capture. Since the
        voltage of a triangle is uniformly distributed, the sorted scope samples are a straight line
        against their rank; fitting the central 80\,\% and extrapolating to the ends averages the
        \SI{9.9}{mV} scope quantization and the noise.
  \item \textbf{Slow slope.} At \SI{\RampFreq}{Hz} the ramp moves about one LSB every 60 samples,
        so every code is hit many times per crossing and noise only smooths the transitions.
\end{itemize}
The accuracy of $T_k$ is limited by the linearity of the DDS ramp and by the gain and offset
accuracy of the scope; the DNL and INL, which are relative, depend only on the ramp linearity.

\subsection{Dynamic performance: sine fit}
\paragraph{Coherent sampling.} The sine frequencies are $f_{in} = J f_s / M$ with $J$ an odd
prime: the record holds exactly $J$ periods, and since $J$ and $M = 2^{16}$ share no factor every
sample falls on a different phase of the sine, so the record exercises the whole transfer
characteristic instead of repeating a few points. The integer number of periods also places the
fundamental and its harmonics exactly on FFT bins.

\paragraph{Four-parameter fit.} Following IEEE~1057, the record $x_n$ is fitted by least squares to
\begin{equation}
y_n = A\cos(2\pi f t_n + \varphi) + C, \qquad t_n = n / f_s,
\end{equation}
where $A$, $\varphi$ and $C$ are linear parameters for a given $f$, and $f$ is refined around the
nominal value by minimizing the residual. Everything that is not the fitted sine (quantization,
nonlinearity, noise, jitter) stays in the residual, so
\begin{equation}
\text{SINAD} = 20\log_{10}\frac{A/\sqrt{2}}{\sqrt{\tfrac{1}{M}\sum_n (x_n - y_n)^2}}.
\end{equation}
The fitted amplitude $A$ against frequency also gives the amplitude response of the front end
and the ADC sampler; the amplitude measured on the scope at each frequency checks that the
stimulus itself stays flat.

\paragraph{Static model.} To separate static from dynamic errors, the fitted sine is converted back
to volts and quantized with the transition levels $T_k$ measured by the ramp, without noise. The
SINAD of this synthetic record (SINAD$_{st}$) is what the static transfer alone allows; the
difference to the measured SINAD is caused by noise and by frequency dependent effects of the
front end and the sampler.

\subsection{Spectrum, harmonics and folding}
The spectrum uses a 4-term Blackman--Harris window, which keeps the energy of each tone in a few
bins even if the ADC and DDS clocks are slightly off; the power of a tone is summed over $\pm 4$
bins. Sampled at $f_s$, a harmonic $h f_{in}$ above the Nyquist frequency $f_s/2$ appears at its
alias
\begin{equation}
f_h = \left| h f_{in} - f_s \cdot \operatorname{round}\!\left(\frac{h f_{in}}{f_s}\right) \right|,
\end{equation}
so at high input frequencies the harmonics fold back below the fundamental. With $P_1$ the power of
the fundamental, $P_h$ of harmonic $h$, $P_{noise}$ the rest of the band and $P_{spur}$ the largest
other tone:
\begin{equation}
\text{THD} = 10\log_{10}\frac{\sum_{h=2}^{7} P_h}{P_1}, \quad
\text{SNR} = 10\log_{10}\frac{P_1}{P_{noise}}, \quad
\text{SFDR} = 10\log_{10}\frac{P_1}{P_{spur}}.
\end{equation}
Harmonics are given in dBc, relative to the fundamental; the spectrum plots are in dBFS, relative to
a full-scale sine of $2^{N-1}$ LSB amplitude.

\subsection{What the differential front end can and cannot cause}
A balanced differential path cancels even-order distortion: a nonlinearity that is the same on both
legs adds the same even terms to $v_p$ and $v_n$, and they vanish in $v_p - v_n$. A gain or phase
imbalance between the legs of the single-ended to differential conversion breaks the cancellation,
so a rising second harmonic points to the front end. Bandwidth limits in the front end or in the
sampling network of the ADC reduce the fitted amplitude at high frequency while the input stays
flat. Errors tied to specific codes, such as narrow or wide codes at the transitions of one bit,
cannot come from an analog front end: they are mismatches of the capacitive DAC of the SAR ADC,
which also produce mostly odd harmonics for a sine centered in the range.
